\documentclass[
	paper=a4,	%Seitengröße
	11pt,		%Schriftgröße
	DIV=11,		%Satzspiegel
	%twoside, %nur für drucken
	openright, ]{scrreprt}%immer rechts kapitel beginnen

%Default-einstellungen für das Dokument

\usepackage[T1]{fontenc}	%Schriftcodierung
\usepackage[utf8]{inputenc}	%Dateicodierung
\usepackage[ngerman]{babel} %Sprachpaket(Silbentrennung,überschriften, Abbildungen)
\usepackage{csquotes}
%Literatur anpassen
\usepackage[ natbib=true, style=numeric,sorting=none]{biblatex}
\addbibresource{ref/ref_liste.bib}

%Schriftanpassungen
\usepackage{lmodern} %Vektorschrift
\renewcommand{\familydefault}{\sfdefault} %serifenlose Schrift
\usepackage{sansmath}

\usepackage[
	automark, %automatische Kolumne
	headsepline %Trennlinie in Kopfzeile
]{scrlayer-scrpage}

\cfoot{}

%wichtige Pakete
\usepackage{graphicx}
\graphicspath{{bilder/}{}}  %Pfad zu den Bildern. mehrere pfade mit {{}{}{}} 

\title{Semantische Segmentierung der Umgebung auf Basis von 3D-Daten}

\author{{\Large Simon Kuhn}}
\date{Wissenschaftliche Arbeit im Zuge des Fachwissenschaftlichen Seminares}
\usepackage{blindtext}

\publishers{
	\begin{tabular}{rl} %jedes C für eine spalte c für centrered

		Erstprüfer: & Prof. Dr. Christian Pfitzner   \\
		Betreuer:   & Prof. Dr. Christian Pfitzner \\ 
		Ausgabedatum: & 23.03.2023 \\
		Abgabedatum: & 31.08.2023  \\

	\end{tabular}
}

\titlehead{	
	\includegraphics[width=0.5\textwidth]{thi_FEI_wb_RGB}
	\hfil
}

\lowertitleback{
	Simon Kuhn
	Robotikstudent
	geb. am 06.10.2001
}

%%%%%%%DOKUMENT%%%%%%%%%%%%%
\begin{document}
\maketitle
\tableofcontents
\chapter{Einleitung}
\section{Hintergrund und Motivation}

In den letzten Jahren hat die Forschung im Bereich der autonomen Fahrzeuge und
der Robotik enorme Fortschritte gemacht. Ein wichtiger Faktor für die
Entwicklung dieser Technologien ist die Fähigkeit, die Umgebung ausreichend
genau zu erkennen und zu verstehen. In diesem Zusammenhang hat die semantische
Segmentierung der Umgebung auf Basis von 3D-Daten eine immer größere Bedeutung
erlangt \cite{.4242023b}. 

Die semantische Segmentierung ist ein Verfahren zur automatischen
Klassifizierung von Objekten und Strukturen in der Umgebung. Dabei werden jedem
Pixel oder jedem Voxel in einem 3D-Modell eine bestimmte semantische Bedeutung
zugeordnet, z.B. Straße, Gebäude, Bäume oder Fahrzeuge. Eine präzise
semantische Segmentierung ist eine wesentliche Voraussetzung für eine
zuverlässige Navigation von autonomen Fahrzeugen und Robotern.
In dieser Arbeit wird die semantische Segmentierung der Umgebung auf Basis von
3D-Daten untersucht. Dabei sollen verschiedene Methoden zur semantischen
Segmentierung untersucht und bewertet werden. Ziel ist es, die Genauigkeit und
Zuverlässigkeit der semantischen Segmentierung zu verbessern und die
Anwendungsmöglichkeiten in der Praxis zu erweitern.

\section{Problemstellung}

Obwohl die semantische Segmentierung der Umgebung auf Basis von 3D-Daten in den
letzten Jahren deutliche Fortschritte gemacht hat, gibt es noch viele offene
Fragen und Herausforderungen. Zum Beispiel sind viele der bestehenden Methoden
für die semantische Segmentierung nur bedingt skalierbar und erfordern viel
Rechenleistung. Zudem sind sie oft sehr empfindlich gegenüber Veränderungen in
der Umgebung, wie z.B. Änderungen der Lichtverhältnisse oder der Perspektive.
Um diese Herausforderungen zu meistern, ist es notwendig, neue Methoden und
Technologien zu entwickeln, die eine robuste und skalierbare semantische
Segmentierung ermöglichen. Diese Arbeit trägt dazu bei, indem sie verschiedene
Methoden und Technologien für die semantische Segmentierung der Umgebung auf
Basis von 3D-Daten untersucht und bewertet.



\chapter{Sensoren zur Erfassung von 3D-Daten}

\section{LiDAR-Sensoren}
LiDAR-Sensoren, die auch unter dem Namen Light Detection and Ranging-Sensoren
bekannt sind, stellen eine weit verbreitete Technologie zur Erfassung von
3D-Daten dar. Sie basieren auf dem Einsatz von Laserstrahlen, welche
ausgesendet werden und von Objekten in der Umgebung reflektiert werden. Dabei
kann zwischen Time of Flight (TOF) LiDAR und phasenbasiertem LiDAR
unterschieden werden. Während TOF-LiDAR die Distanz über eine Messung der
Laufzeit der lichtwelle bestimmt, erfolgt die Entfernungsmessung beim
phasenbasierten LiDAR über die Auswertung der Phasenverschiebung der vom Objekt
reflektierten Lichtwelle. Hierdurch können LiDAR-Sensoren hochgenaue
Entfernungen zu den reflektierenden Objekten erfassen aus denen sich
detaillierte 3D-Punktwolken erzeugen lassen, welche die Geometrie und räumliche
Verteilung von Objekten in der Umgebung darstellen. Zusätzlich lassen sich
LiDAR-Sensoren in Scanning-LiDAR und Non-Scanning-LiDAR untergliedern.
Non-Scanning-LiDAR nutzt dabei einen statischen Laserstrahl, während
Scanning-LiDAR einen sich bewegenden Laserstrahl nutzt. \cite{8529992}

\section{Tiefenkameras}
Tiefenkameras basieren auf verschiedene Funktionsweisen, um Entfernungen zu
messen. Im Bereich der semantische Segmentierung kommen besonders Sensoren auf
Basis von Stereo-Vision, Time-of-Flight, oder Structured Light zum Einsatz.
Kamerasysteme, die auf dem Prinzip der Stereo-Vision basieren, werden als
Stereo Kameras bezeichnet. Bei diesen werden zwei räumlich getrennte Kameras
verwenden, die gemeinsam Bilder von derselben Szene aus zwei leicht
unterschiedlichen Perspektiven aufnehmen. Der dabei entstehende horizontale
Versatz der beiden Bilder wird als Disparität bezeichnet. Aus diesem lassen
sich Tiefeninformationen des betrachteten Objektes berechnen. Bei Time of
Flight Kameras wird ein Lichtpuls im Infrarotbereich ausgesendet und von
Objekten in der Umgebung reflektiert. Über die Laufzeit der Infrarotwelle lässt
sich die Entfernung berechnen. 3D-Kameras auf Basis von Structured Light
projizieren ein Muster auf das zu betrachtende Objekt. Aus der Verzerrung des
dessen, lassen sich Tiefeninformationen berechnen. Die meisten Tiefenkameras
stellen dabei die Tiefeninformationen in einem Bild aus Graustufen dar.
Zusätzlich gibt es auch RGB-D Kameras, welche neben einer Strucutred Light oder
TOF-Kamera auch über eine RGB-Kamera verfügen. Diese haben häufig eine höhere
räumliche Auflösung und sind in der Lage zusätzlich Farbinformationen
aufzunehmen, besitzen jedoch eine deutlich kleineren Arbeitsbereich.

\section{Passive und aktive Sensoren}
Aktive und passive Sensoren sind zwei grundlegende Arten von Sensoren, die in
der 3D-Bildgebung eingesetzt werden. Aktive Sensoren, wie beispielsweise
LIDAR-Sensoren, senden eine Energiequelle aus, wie zum Beispiel Laser oder
Infrarot, um Informationen über das Objekt zu sammeln. Die reflektierten
Signale werden von der Sensor-Einheit aufgenommen und zur Berechnung von
Tiefeninformationen verwendet. Im Gegensatz dazu erfordern passive Sensoren,
wie Tiefenkameras, keine aktive Energiequelle, sondern nutzen das natürliche
Licht, das von der Umgebung reflektiert wird. Passive Sensoren messen das
reflektierte Licht, um die Tiefeninformationen zu berechnen. Obwohl beide
Sensortypen Tiefeninformationen sammeln können, haben sie unterschiedliche
Stärken und Schwächen. LIDAR-Sensoren sind beispielsweise oft präziser, aber
teurer als Tiefenkameras, die einfacher zu verwenden und kosteneffektiver sind.

\section{Auswahl von Sensoren für die semantische Segmentierung}

Die Auswahl der geeigneten Sensoren für die semantische Segmentierung hängt von
verschiedenen Faktoren ab, wie den Anforderungen der spezifischen Anwendung,
den Umgebungsbedingungen, dem Budget und den gewünschten Ergebnissen. In diesem
Abschnitt werden Kriterien und Überlegungen zur Auswahl von Sensoren für die
semantische Segmentierung auf Basis von 3D-Daten diskutiert.

Anforderungen der Anwendung: Die Anforderungen der spezifischen Anwendung, in
der die semantische Segmentierung durchgeführt werden soll, sind entscheidend
für die Auswahl der geeigneten Sensoren. Hierbei können Aspekte wie die
benötigte Genauigkeit, die räumliche Auflösung, die Reichweite, die
Echtzeitfähigkeit und die Umgebungsbedingungen eine Rolle spielen. Zum Beispiel
erfordern Anwendungen im Bereich der autonomen Fahrzeuge möglicherweise
Sensoren mit hoher Reichweite und Genauigkeit, während Anwendungen im
Innenbereich möglicherweise Sensoren mit höherer räumlicher Auflösung und
Echtzeitfähigkeit benötigen.

Umgebungsbedingungen: Die Umgebungsbedingungen, in denen die semantische
Segmentierung durchgeführt werden soll, können die Auswahl der Sensoren
beeinflussen. Beispielsweise können schlechte Beleuchtungsbedingungen, wie
Dunkelheit oder Blendung durch Sonnenlicht, die Leistung von passiven Sensoren
wie Kameras beeinträchtigen, während Lidar-Sensoren unabhängig von den
Beleuchtungsbedingungen arbeiten können. Ebenso können Umgebungen mit vielen
Hindernissen oder komplexen Geometrien die Leistung von Sensoren beeinflussen
und die Wahl von geeigneten Sensoren beeinflussen.

Budget: Das Budget ist ein wichtiger Faktor bei der Auswahl von Sensoren.
Unterschiedliche Sensoren können unterschiedliche Kosten haben, sowohl in der
Anschaffung als auch in der Wartung. Lidar-Sensoren sind in der Regel teurer
als Kameras oder Tiefenkameras, aber auch die Kosten für diese Sensoren haben
sich in den letzten Jahren reduziert. Bei der Auswahl von Sensoren ist es
wichtig, das Budget im Auge zu behalten und eine sorgfältige
Kosten-Nutzen-Analyse durchzuführen.

Gewünschte Ergebnisse: Die gewünschten Ergebnisse der semantischen
Segmentierung können ebenfalls die Auswahl der Sensoren beeinflussen. Je
nachdem, welche Arten von Objekten oder Strukturen in der Umgebung segmentiert
werden sollen, können bestimmte Sensoren besser geeignet sein als andere. Zum
Beispiel können Lidar-Sensoren aufgrund ihrer präzisen Tiefeninformationen und
Reichweite gut geeignet sein, um Objekte wie Straßen, Gebäude oder Bäume zu
segmentieren, während Kameras oder Tiefenkameras besser für die Segmentierung
von Fußgängern oder Fahrzeugen geeignet sein können.
\input{kapitel/datengrundlage_und_Vorverarbeitung} 
\chapter{Grundlegende Verfahren der semantischen Segmentierung}

Ziel der Semantischen Segmentierung von 3D-Daten ist es, jedem Punkt im Raum
einer bestimmten Kategorie zuzuordnen und dadurch Bereiche des Bildes in
klassifizierte Objekte zu unterteilen. Hierfür gibt es verschiedene Ansätze,
die auf erweiterten neuronalen Netzen basieren.

\subsection{Convolutional Neural Networks (CNNs)}
Convolutional Neural Networks (CNNs) sind eine erweiterte Art von künstlichen
neuronalen Netzen, die speziell für die Verarbeitung von Bildern entwickelt
wurden. Sie bestehen aus mehreren Schichten, darunter Convolutional Layers,
Pooling Layers und Fully Connected Layers, die miteinander verbunden sind. Die
Architektur der CNNs ist dabei nicht vorgegeben, folgt aber in der Praxis immer
einer ähnlichen Vorlage. Ein Input-Layer beinhaltet die Pixelwerte des Bildes.
Darauf folgen in der Regel ein oder zwei Convolutional-Layers, woraufhin sich
ein Pooling-Layer anreiht. Die Kombination aus Convolutional- und Pooling-Layer
kann dabei je nach Komplexität beliebig oft im Netz vorkommen. Am Ende folgt
ein Fully-Connected-Layer an den der Output anknüpft. Dieser gibt das
Klassifizierungsergebnis des betrachteten Bildes durch das Netzwerk aus. Die
Convolutional-Layers können vereinfacht als Filter-Layer betrachtet werden und
extrahieren markante Merkmale aus den Eingabebildern, um die Klassifikation zu
erleichtern. Dabei wird jeder Bereich des Bildes pro Layer mit einem Kernel
gefaltet und erzeugen eine 2D-Aktivierungskarte. Die kernels besitzen dabei
häufig kleine räumliche Dimensionalitäten, erstrecken sich aber über die
Gesamte Tiefe des Eingangsbildes. Die Anschließend folgenden Pooling-Layers
dienen dazu, die Dimensionen der Ausgabe des Convolutional-Layers zu reduzieren
und somit die Rechnenkomplexität des Modells zu verringern. Besonders häufig
kommen dabei sogenannte Max-Pooling-Layers zum Einsatz. Bei diesen wird ein
Kernel einer beliebigen Größe, ohne zu überlappen, über die Aktivierungskarte
geschoben. Dabei wird nur der größte Aktivierungswert eines Fensterbereiches
übernommen und so die Dimensionalität der Aktivierungskarte stark reduziert.
Die Fully Connected Layers am Ende des Netzes verarbeiten schließlich die
extrahierten Aktivierungen und versuchen daraus Klassifizierungsergebnisse zu
gewinnen. Der Hauptanwendungsbereich von CNNs liegt dabei in der Klassifikation
von Bildern in vorbestimmte Kategorien. \cite{11262015}.

\begin{figure}
    \centering
    \includesvg[width=0.8\textwidth]{bilder/cnn_own.svg}
    \caption{Aufbau eines Convolutional Neural Networks}
    \label{fig:meinbild}
\end{figure}

\subsection{Fully Convolutional Networks (FCNs)}
Fully Convolutional Networks (FCNs) sind eine Weiterentwicklung von CNNs, die
speziell für die Aufgabe der semantischen Segmentierung von Bildern entwickelt
wurden. Im Gegensatz zu herkömmlichen CNNs, die für die Klassifizierung und
Erkennung von Objekten in Bildern ausgelegt sind, können FCNs jedes Pixel eines
Eingabebildes klassifizieren und somit die räumlichen Informationen der
klassifizierten Bilder beibehalten. FCNs verwenden dabei ausschließlich
Convolutional-Layers und Pooling-Layers, nicht aber Fully-Connected-Layers.
Dies ermöglicht es eine Merkmalskarte des Eingabebildes zu erzeugen, auf der
jedes Pixel einer bestimmten Klasse zugeordnet wird. Durch die Verwendung von
FCNs können somit komplizierte Zusammenhänge innerhalb von Bildern auf der
Ebene der Pixel identifiziert werden, was für Anwendungen wie die autonome
Navigation oder Objekterkennung von großer Bedeutung ist. \cite{7803544}
Bekannte und erfolgreiche Verfahren wie das U-Net, verwenden dabei häufig eine
Encode-Decoder-Architektur, um das Klassifizierungsergebnis in der
Dimensionalität des Eingangsbildes darzustellen.

\subsection{Encoder-Decoder-Architekturen}
Encoder-Decoder-Architekturen stellen eine spezielle Architektur von neuronalen
Netzen dar. Der Name leitet sich aus deren Aufbau ab, welcher aus zwei
Hauptkomponenten, einem Encoder und einem Decoder, besteht. Der Encoder
verwendet typischerweise Convolutional und Pooling-Layers, um das Eingabebild
schrittweise in eine kompakte, abstrakte Repräsentation zu komprimieren, die
die Merkmale des Eingabebildes stark reduziert enthält. Dabei werden die
Positionen der maximalen Aktivierungen der Ebene während des Max-Pooling
Prozesses gespeichert und als Pooling-Indizes bezeichnet. Der Decoder verwendet
häufig Deconvolutional-Neuronale-Networks, um das Ergebnis des
Encoder-Netzwerkes wieder in die Dimensionalitäten des Eingangsbildes
zurückzuführen. Das Upsampling erfolgt dabei nicht linear, sondern verwendet
die Pooling-Indizes des zugehörigen Encoding-Schrittes, um die Aktivierungen an
der richtigen Position des Bildes wiederherzustellen. Dies geschieht über
sogenannte Skip-Verbindungen. Die Semantische Segmentierung erfolgt dabei am
Ende des Decoding Netzwerkes. Dabei wird die Merkmalskarte des Decoders in eine
Wahrscheinlichkeitsverteilung bezüglich der zu klassifizierenden Klassen
umgewandelt. Die geschieht häufig über die Softmax-Funktion. \cite{7803544}

\begin{figure}
    \centering
    \includesvg[width=0.8\textwidth]{bilder/FCN_own_new.svg}
    \caption{Encoder Decoder Architektur ähnlich U-Net}
    \label{fig:meinbild}
\end{figure}

\subsection{Region-based Convolutional Neural Networks (R-CNNs)}
Region-based Convolutional Neural Networks (R-CNNs) sind eine Weiterentwicklung
von Convolutional Neural Networks. Im Gegensatz zu herkömmlichen CNNs, die eine
feste Größe der Eingabebilder erfordern, verwenden R-CNNs eine Region Proposal
Technik, um Regions of Interest (ROI) innerhalb des Bildes zu detektieren.
Anschließend wird nur auf die ausgewählten Bereiche eine Klassifizierung
durchgeführt. R-CNNs erzielen dadurch eine höhere Genauigkeit als herkömmliche
CNNs bei der Erkennung von Objekten in Bildern und werden daher häufig in der
Robotik und im autonomen Fahren eingesetzt. Dies liegt daran, dass sie weniger
Rechenleistung benötigen und gleichzeitig Störeinflüsse, die außerhalb der ROI
liegen, keinen negativen Einfluss auf die Klassifizierung nehmen können. \cite{8237548}
\chapter{Anwendungszenarien der semantischen Segmentierung}
\section{Autonomes Fahren}
Eines der wichtigsten Einsatzgebiete der semantischen Segmentierung ist der
Bereich des autonomen Fahrens. Hierbei ist eine präzise und zuverlässige
Wahrnehmung der Umgebung unumgänglich. Dabei ist eine der größten
Herausforderungen die Verarbeitung der großen Datenmengen, die von den Sensoren
erfasst werden. Die Verarbeitungsgeschwindigkeit muss dabei extrem hoch sein,
um in Echtzeit auf die Umgebung reagieren zu können. Durch die semantische
Segmentierung ist es dem Fahrzeug möglich, wichtige Informationen über
Straßenverhältnisse, Verkehrszeichen, Fußgänger, Fahrzeuge und andere
Hindernisse zu sammeln. Mithilfe einer präzise nKlassifizierung jedes Pixels im
Bild ist es dem Fahrzeug möglich, Objekte in der Umgebung kontextbezogen zu
erkennen und situationsgerecht zu reagieren.\cite{9000872}

\section{Robotik in der Industrie}

Die semantische Segmentierung spielt eine wichtige Rolle in der
Industrierobotik, insbesondere wenn es darum geht, stark variierende Produkte
bezüglich ihres Formfaktors zu handhaben. Eine Herausforderung besteht darin,
die Produkte zu erkennen und zu klassifizieren. Dies ermöglicht es dem Roboter
zwischen unterschiedlichen Teilen zu unterscheiden und bauteilspezifisch zu
reagieren. So kann beispielsweise eine optimale Griffposition des Bauteils
ermittelt werden, oder Objekten beim Anfahren gezielt ausgewichen werden. Auch
eine Manipulation komplexer Objekte kann dadurch ermöglicht werden. \cite{liu2020real}

\section{Augmented Reality}
Einsatzgebiete der semantischen Segmentierung im Bereich von Augmented Reality
(AR) sind vielfältig. Diese kann beispielsweise in der industriellen Fertigung
eingesetzt werden, um komplexe Montage- und Wartungsprozesse zu unterstützen.
Durch semantische Segmentierung können AR-Systeme bereits platzierte Teile und
Komponenten identifizieren, Fehler erkennen und den Arbeitern visuelle
Montageanleitungen in Echtzeit bereitstellen. Dies erhöht die Effizienz,
Präzision und Sicherheit bei der Durchführung von Aufgaben. In der Architektur
und im Bauwesen ermöglicht die semantische Segmentierung in Kombination mit AR
eine Visualisierung von Gebäuden und Infrastrukturen. Architekten, Designer und
Ingenieure können virtuelle Modelle in die reale Umgebung einfügen, um Entwürfe
auf Kollisionen zu überprüfen. Die semantische Segmentierung ist dabei eine
entscheidende Komponente, um die Umgebung zu analysieren, zu verstehen und
interaktiv nutzbar zu machen. Dies ermöglicht es AR-Systemen, die Umgebung zu
analysieren und virtuelle Inhalte realitätsnah zu platzieren und mit diesen zu
interagieren. \cite{nee2012augmented,ko2020novel}

\section{Landwirtschaft}

Auch in der Landwirtschaft hat die semantische Segmentierung große Potenziale.
Durch die präzise Klassifizierung von verschiedenen Pflanzenarten und
Strukturen auf landwirtschaftlichen Flächen, können Landwirte wichtige
Informationen gewinnen, um ihre Produktivität und Effizienz zu steigern. Mit
Hilfe der semantischen Segmentierung können Unkraut und Nutzpflanzen
unterschieden werden. Dies ermöglicht eine gezielte Anwendung von Düngemitteln
und Pestiziden, sowie einer effizienten Bewässerung. Darüber hinaus kann die
semantische Segmentierung auch bei der Erkennung von Schädlingen und
Krankheiten helfen, um rechtzeitig Gegenmaßnahmen zu ergreifen. Durch die
genaue Analyse und Segmentierung landwirtschaftlicher Flächen bietet die
semantische Segmentierung ein wertvolles Werkzeug, um die Landwirtschaft
nachhaltiger, produktiver und umweltfreundlicher zu gestalten. \cite{8460962}

\section{Medizin}
Semantische Segmentierung spielt eine entscheidende Rolle im Bereich der
Medizintechnik und hat das Potenzial, unterstützend bei medizinische Diagnosen
und Behandlungen zu wirken. Durch die präzise Segmentierung von anatomischen
Strukturen und Geweberegionen in medizinischen Bildern wie CT-Scans oder
MRT-Aufnahmen können Ärzte einen besseren Einblick in die Anatomie des
Patienten erhalten. Beispielsweise können Organe für einen besseren Überblick
automatisch eingefärbt werden, oder verdächtige Stellen in der Aufnahme
markiert werden. So kann der untersuchende Arzt beispielsweise Verletzungen
oder potenzielle Tumore gezielt betrachten. Diese detaillierte Analyse
ermöglicht eine präzisere Diagnosestellung und kann bei der Planung und
Durchführung von chirurgischen Eingriffen, sowie der Überwachung des
Fortschreitens von Krankheiten eingesetzt werden. Die semantische Segmentierung
erleichtert auch die Erstellung von 3D-Modellen anatomischer Strukturen, die
für die präoperative Planung und Simulation von komplexen Operationen verwendet
werden können. Dabei können beispielsweise markante Stellen einer CT-Aufnahme
bei einer endoskopischen Operation wieder erkannt werden und dem operierenden
Arzt unterstützend eingeblendet werden. Dies kann unter anderem auch mithilfe
von AR geschehen. \cite{asgari2021deep,tanzi2021real}
\chapter{State-of-the-Art Verfahren zur semantischen Segmentierung von 3D-Daten}
\subsection{PointNet}
PointNet++ ist eine fortschrittliche Methode zur Verarbeitung von Punktwolken
in der 3D-Bildverarbeitung, die speziell für die Aufgaben der Klassifikation,
Segmentierung und Erkennung von Objekten entwickelt wurde. Es basiert auf der
PointNet-Architektur und erweitert sie durch die Integration einer
hierarchischen Struktur, um lokale und globale Kontextinformationen in
Punktwolken zu erfassen.

Die Funktionsweise von PointNet++ besteht aus mehreren aufeinanderfolgenden
Schritten. Zunächst wird die ursprüngliche Punktwolke in eine hierarchische
Struktur unterteilt. Dies wird erreicht, indem die Punktwolke in immer kleinere
Unterteilungen aufgeteilt wird, wobei auf jeder Hierarchieebene ein separater
PointNet-Block angewendet wird.

Auf jeder Hierarchieebene führt der PointNet-Block zwei grundlegende
Operationen durch: Das Pooling und die Verarbeitung von Punktwolken. Das
Pooling dient dazu, die relevanten Merkmale aus den Punkten auf jeder
Hierarchieebene zu extrahieren. Dabei werden aggregierte Merkmale auf einer
höheren Ebene erzeugt, um den globalen Kontext der Punktwolke zu erfassen.

Die Verarbeitung von Punktwolken auf jeder Hierarchieebene beinhaltet die
Anwendung von PointNet auf die Punkte innerhalb der Teilbereiche der
Punktwolke. Dies ermöglicht die Extraktion von lokalen Merkmalen, die
spezifisch für bestimmte Regionen oder Strukturen in der Punktwolke sind.

Die hierarchische Struktur von PointNet++ ermöglicht es, sowohl lokale als auch
globale Kontextinformationen zu erfassen und zu integrieren. Durch die
Verwendung mehrerer PointNet-Blöcke auf verschiedenen Hierarchieebenen können
feinere Details auf lokaler Ebene berücksichtigt werden, während gleichzeitig
der globale Kontext der Punktwolke erhalten bleibt.

Die Ausgabe von PointNet++ ist eine repräsentative Merkmalsdarstellung, die die
Informationen über die semantische Struktur der Punktwolke enthält. Diese
Merkmale können dann für verschiedene Aufgaben wie Klassifikation,
Segmentierung oder Erkennung von Objekten verwendet werden.

Insgesamt ermöglicht die Funktionsweise von PointNet++ die effektive
Verarbeitung von Punktwolken und die Erfassung von lokalen und globalen
Kontextinformationen. Es hat sich als leistungsstarkes Verfahren erwiesen, um
komplexe Strukturen und Muster in 3D-Daten zu erfassen und die Genauigkeit und
Zuverlässigkeit bei der Klassifikation und Segmentierung von Punktwolken zu
verbessern.
\subsection{3D U-Net}

3D U-Net ist ein leistungsstarkes Verfahren für die semantische Segmentierung von 3D-Daten. Es basiert auf der beliebten 2D U-Net-Architektur, die für die Bildsegmentierung entwickelt wurde, und wurde speziell für die Verarbeitung von Volumendaten angepasst.

Die Funktionsweise von 3D U-Net kann in zwei Hauptphasen unterteilt werden: den
Encoder-Teil und den Decoder-Teil.

Im Encoder-Teil werden die Eingabedaten schrittweise abwärts durch das Netzwerk
geleitet, um Merkmale auf verschiedenen Abstraktionsebenen zu extrahieren. Dies
erfolgt durch eine Kombination aus 3D-Convolutional-Layern und
Pooling-Operationen, die die räumliche Auflösung der Daten reduzieren. Dabei
werden lokalisierte Merkmale erfasst, um wichtige Informationen über die
Struktur und den Kontext der 3D-Daten zu gewinnen.

Der Decoder-Teil arbeitet in umgekehrter Richtung und verwendet
Upsampling-Operationen, um die räumliche Auflösung schrittweise zu erhöhen.
Dabei werden die Merkmale aus dem Encoder-Teil mit den entsprechenden Merkmalen
auf höheren Auflösungsebenen kombiniert, um detailliertere und präzisere
Segmentierungsergebnisse zu erzielen. Dieser Prozess wird durch sogenannte
"Skip Connections" ermöglicht, die es dem Decoder ermöglichen, sowohl lokale
als auch globale Informationen zu nutzen.

Während des Trainings wird das 3D U-Net-Modell mit gelabelten Trainingsdaten
trainiert, um die Gewichtungen der Netzwerkparameter zu optimieren. Dies
geschieht durch den Vergleich der vorhergesagten Segmentierungsergebnisse mit
den tatsächlichen Labels. Durch den Einsatz von Verlustfunktionen wie der
Kreuzentropie wird das Modell kontinuierlich verbessert und kann genaue
semantische Segmentierungen von neuen, nicht-gelabelten Daten vorhersagen.

Die Funktionsweise von 3D U-Net hat sich in verschiedenen Anwendungen bewährt,
darunter medizinische Bildgebung, Robotik, geografische Kartierung und mehr. Es
ermöglicht eine präzise Segmentierung von 3D-Daten und eröffnet Möglichkeiten
zur Analyse und Interpretation komplexer Strukturen und Merkmale in
volumetrischen Daten.

\subsection{OctNet}

OctNet ist ein leistungsstarkes Verfahren zur Verarbeitung und Segmentierung
von 3D-Daten, das auf der Verwendung von Octrees basiert. Octrees sind eine
hierarchische Datenstruktur, die es ermöglicht, den Raum in kleine Voxelblöcke
zu unterteilen und gleichzeitig eine adaptive Auflösung bereitzustellen.

Die Funktionsweise von OctNet beruht auf der effizienten Organisation und
Verarbeitung von 3D-Daten. Zunächst wird die 3D-Datenrepräsentation in ein
Octree umgewandelt, wobei jeder innere Knoten des Baums acht Unterknoten hat,
die den Raum in kleinere Voxelblöcke unterteilen. Dadurch können Bereiche mit
höherer Detailgenauigkeit mehr Unterknoten haben, während weniger detaillierte
Bereiche weniger Unterknoten aufweisen.

Das OctNet-Modell besteht aus einer Kombination von Convolutional Neural
Networks (CNNs) und speziellen Operationen, die auf die Struktur des Octrees
abgestimmt sind. Durch die Verwendung von Convolutional-Operationen auf den
Octree-Daten werden Merkmale auf verschiedenen Auflösungsebenen erfasst und
repräsentiert. Dies ermöglicht eine effiziente und adaptive Verarbeitung von
3D-Daten, da die Operationen nur auf den relevanten Voxelblöcken des Octrees
durchgeführt werden.

Während des Trainings wird das OctNet-Modell mit gelabelten Trainingsdaten
trainiert, um die Gewichtungen der Netzwerkparameter zu optimieren. Dabei wird
die Beziehung zwischen den Eingabedaten und den entsprechenden
Segmentierungskarten erlernt. Durch die Anpassung der Gewichtungen können
genaue Segmentierungsergebnisse erzielt werden, die die semantischen Strukturen
in den 3D-Daten korrekt erfassen.

Die Funktionsweise von OctNet ermöglicht eine effiziente Verarbeitung und
Segmentierung von 3D-Daten mit variabler Auflösung. Es hat sich in
verschiedenen Anwendungen wie der Segmentierung von Punktwolken, der Analyse
von 3D-Modellen und der medizinischen Bildgebung als wirksam erwiesen. Durch
die Verwendung von Octrees bietet OctNet eine leistungsstarke Methode zur
Erfassung und Analyse von komplexen Strukturen in 3D-Daten.
\input{kapitel/Herausforderungen und zukünftige Entwicklungen}
\chapter{Zusammenfassung}
In dieser Arbeit wurden die Grundlagen und Anwendungen der semantischen
Segmentierung auf Basis von 3D-Daten dargestellt. Es wurde auf die
Funktionsweise verschiedene Sensoren zu Gewinnung von 3D-Daten eingegangen und
deren Vor- und Nachteile erläutert. Daraufhin wurde auf die Datengrundlagen
eingegangen, auf denen verschiedene Modelle für die Semantische Segmentierung
basieren. Auch mögliche Vorverarbeitungsschritte für eine effizientere und
präzisere Segmentierung wurden erläutert. Dabei ging es vor allem um die
Filterung, Normalenberechnung und Möglichkeiten für das Downsamplings eines
3D-Datensatzes. Anschließend wurde auf die Grundlagen der semantischen
Segmentierungsverfahren eingegangen. Besonders Convolutional Neural Networks
(CNNs), Fully Convolutional Networks (FCNs), Region-based Convolutional Neural
Networks (R-CNNs) und Encoder-Decoder-Architekturen wurden erläutert. Dabei hat
jedes Verfahren spezifische Einsatzgebiete und Vor- und Nachteile. Dabei wurden
zwei verbreitete Verfahren für die semantische Segmentierung aufgeführt und
kurz erläutert. Um die Wichtigkeit der semantischen Segmentierung zu
verdeutlichen, wurden anschließend einige Anwendungsbereiche der semantischen
segmentieren genannt und deren Bedeutung in diesem Bereich verdeutlicht. Am
Ende der Arbeit gab es einen Ausblick auf die bestehenden Herausforderungen und zukünftig
zu erwartende Entwicklungen im Bereich der semantischen Segmentierung.

\printbibliography
\end{document}